\section{Metodologia}

\subsection{Dataset sintetico}

Il dataset viene generato campionando casualmente i parametri dell'opzione e del mercato.
Ogni osservazione contiene:

\begin{equation}
    x = (S_0, K, T, r, \sigma),
\end{equation}

e il target e':

\begin{equation}
    y = C_{BS}(S_0, K, T, r, \sigma).
\end{equation}

I range iniziali sono:

\begin{table}[H]
    \centering
    \begin{tabular}{lll}
        \toprule
        Variabile & Descrizione & Range \\
        \midrule
        $S_0$ & prezzo iniziale del sottostante & $[50, 150]$ \\
        $K$ & strike & $[50, 150]$ \\
        $T$ & maturity in anni & $[0.1, 2]$ \\
        $r$ & tasso risk-free & $[0, 0.05]$ \\
        $\sigma$ & volatilita' & $[0.1, 0.6]$ \\
        \bottomrule
    \end{tabular}
    \caption{Range usati per la generazione del dataset sintetico.}
\end{table}

La scelta di dati sintetici rende il problema controllato e permette di valutare direttamente l'errore della rete rispetto alla funzione target esatta.

\subsection{Neural network}

Il modello usato e' una rete neurale feed-forward per regressione.
La scelta e' motivata dal fatto che reti feed-forward sufficientemente ampie sono approssimatori universali di funzioni sotto ipotesi generali \cite{cybenko1989approximation,hornik1989multilayer}.
L'architettura iniziale e':

\begin{equation}
    5 \rightarrow 64 \rightarrow 64 \rightarrow 32 \rightarrow 1,
\end{equation}

con funzioni di attivazione ReLU negli strati nascosti.
La rete approssima la funzione:

\begin{equation}
    f_\theta(S_0, K, T, r, \sigma) \approx C_{BS}.
\end{equation}

La loss ottimizzata e' il mean squared error:

\begin{equation}
    \mathcal{L}(\theta) =
    \frac{1}{n}
    \sum_{i=1}^{n}
    \left(f_\theta(x_i) - y_i\right)^2.
\end{equation}

\subsection{Training e validazione}

I dati vengono divisi in training set, validation set e test set.
Gli input vengono standardizzati tramite \texttt{StandardScaler}; anche il target viene scalato durante il training per migliorare stabilita' numerica e convergenza.
La rete viene addestrata con Adam, early stopping e seed fissato per garantire riproducibilita'.

\subsection{Metriche}

Le metriche principali sono:

\begin{itemize}
    \item Mean Absolute Error (MAE);
    \item Root Mean Squared Error (RMSE);
    \item coefficiente di determinazione $R^2$;
    \item MAPE calcolato solo sui prezzi maggiori di 1, per evitare instabilita' quando il prezzo teorico e' vicino a zero.
\end{itemize}
